% 第4章第3节：Sobolev空间（成员条件与范数、弱微分的定义域、完备性与Hilbert结构）。

\section{Sobolev空间}
\label{sec:sob-spaces}

弱导数的存在只回答了输出是什么类型的对象，还没有限制函数和导数有多大。$|x|$在整条实轴上有弱导数$\operatorname{sgn}x$，可两者都不可积；$\sin(nx)$的导数是$n\cos(nx)$，函数大小不变而导数随$n$无界。要把“变化受控”变成可以计算误差、取极限的条件，需要把函数和它的各阶弱导数一起放进范数。本节建立这样的空间，并证明极限过程中导数关系不会丢失。

\subsection{成员条件与Sobolev范数}
\label{sec:sob-norms}

我们希望同时控制函数本身和它直到某一阶的全部弱导数，并以$L^p$范数衡量每一项。

\begin{definition}[Sobolev空间]
\label{def:sob-sobolev-space}
设$\Omega\subset\mathbb R^d$为开集，$k\ge0$为整数，$1\le p<\infty$。称
\[
 W^{k,p}(\Omega)=\bigl\{u\in L^p(\Omega):\ \text{对每个}|\alpha|\le k,\ D^\alpha u\text{ 存在且属于 }L^p(\Omega)\bigr\}
\]
为\term{Sobolev空间}（Sobolev space），其中$D^0u=u$，并赋以\term{Sobolev范数}（Sobolev norm）
\begin{equation}
 \|u\|_{W^{k,p}}=\Bigl(\sum_{|\alpha|\le k}\|D^\alpha u\|_{L^p}^p\Bigr)^{1/p} .
 \label{eq:sob-sobolev-norm}
\end{equation}
\end{definition}

这里的元素仍是几乎处处等价类，各阶弱导数由定理~\ref{thm:sob-weak-derivative-unique}确定到几乎处处。$L^p(\Omega)$的函数在每个紧集上可积（在紧集$K$上由Hölder不等式，$\int_K|u|\le|K|^{1-1/p}\|u\|_{L^p}$），所以定义中的弱导数恒等式都有意义。范数~\eqref{eq:sob-sobolev-norm}是有限个$L^p$范数的$\ell^p$组合：齐次性显然；三角不等式先对每个分量用$L^p$的Minkowski不等式，再对有限维向量$(\|D^\alpha u\|_{L^p})_\alpha$用$\ell^p$的Minkowski不等式；若$\|u\|_{W^{k,p}}=0$，则零阶项$\|u\|_{L^p}=0$，$u=0$。$k=0$时$W^{0,p}=L^p$。

\paragraph{用前面的例子检验成员条件}
在区间$(-1,1)$上，$|x|$和$\operatorname{sgn}x$都属于$L^p$，而$\operatorname{sgn}x$就是$|x|$的弱导数，所以$|x|\in W^{1,p}(-1,1)$对每个$1\le p<\infty$成立。直接计算$\|\,|x|\,\|_{L^p}^p=2/(p+1)$，$\|\operatorname{sgn}\|_{L^p}^p=2$，于是
\[
 \|\,|x|\,\|_{W^{1,p}(-1,1)}=\Bigl(\frac{2}{p+1}+2\Bigr)^{1/p},\qquad p=2\text{ 时等于 }\sqrt{8/3}\approx1.633 .
\]
$|x|$却不在$W^{2,p}(-1,1)$中，因为它的二阶广义导数是$2\delta_0$，没有局部可积表示。阶跃函数属于$L^p(-1,1)$，但没有弱导数，不在任何$W^{1,p}(-1,1)$中；Cantor函数同样如此。

高维的例子可以用幂函数构造。在单位球$B\subset\mathbb R^d$上取$u(x)=|x|^{-a}$，$a>0$。它在原点外光滑，$|\nabla u|=a|x|^{-a-1}$。当$a+1<d$时，$\nabla u$局部可积，并且挖去半径为$r$的小球后分部积分产生的球面项不超过常数乘$r^{-a}\cdot r^{d-1}$，随$r\to0$趋于零，所以经典梯度就是弱梯度。由上一节的幂函数阈值，$u\in L^p(B)$当且仅当$ap<d$，$\nabla u\in L^p(B)$当且仅当$(a+1)p<d$。因此
\begin{equation}
 |x|^{-a}\in W^{1,p}(B)\quad\Longleftrightarrow\quad(a+1)p<d .
 \label{eq:sob-power-membership}
\end{equation}
例如$d=3$、$p=2$时条件是$a<1/2$：$|x|^{-1/4}$属于$W^{1,2}(B)$，却在原点附近无界。与一维的式~\eqref{eq:sob-1d-sup-bound}对照，这已经提示高维Sobolev函数不必有界，第~\ref{sec:sob-embedding}节会回到这个例子。

\paragraph{零阶项的作用}
范数~\eqref{eq:sob-sobolev-norm}中零阶项$\|u\|_{L^p}$不能省略。若只累计导数，例如在$(0,1)$上取$|u|_{W^{1,p}}=\|u'\|_{L^p}$，那么每个非零常数都有$|c|_{W^{1,p}}=0$，这只是一个\term{半范数}（seminorm）：它满足齐次性和三角不等式，却不能区分相差一个常数的函数。导数控制变化，零阶项控制整体的水平；只有在用边界条件或均值条件排除常数以后，导数半范数才可能单独控制函数，这是第~\ref{sec:sob-poincare}节的Poincaré不等式要回答的问题。

\paragraph{弱微分作为有界算子}
在选定的范数下，弱微分是有界线性算子。对$|\alpha|\le k$，把$D^\alpha$看作映射$W^{k,p}(\Omega)\to L^p(\Omega)$。范数中含有$\|D^\alpha u\|_{L^p}^p$这一项，所以
\[
 \|D^\alpha u-D^\alpha v\|_{L^p}=\|D^\alpha(u-v)\|_{L^p}\le\|u-v\|_{W^{k,p}} .
\]
这说明在Sobolev范数下接近的函数，各阶弱导数也在$L^p$中接近。按第~\ref{sec:lfo-boundedness}节的说法，$D^\alpha:W^{k,p}\to L^p$的算子范数不超过$1$。

这并不意味着微分在$L^p$上有界。在$(0,2\pi)$上取$u_n(x)=\sin(nx)$，$\|u_n\|_{L^2}=\sqrt\pi$保持不变，$\|u_n'\|_{L^2}=n\sqrt\pi$却随$n$无界；更何况许多$L^p$函数根本没有弱导数。有界性来自定义域的选择：把输入空间换成$W^{k,p}$，等于预先把导数的大小算进了输入的范数。

\paragraph{本质有界的情形}
$p=\infty$时，$\ell^p$组合改为取最大值。令$W^{k,\infty}(\Omega)$为$u$及所有$|\alpha|\le k$阶弱导数都属于$L^\infty(\Omega)$的函数全体，范数为
\[
 \|u\|_{W^{k,\infty}}=\max_{|\alpha|\le k}\|D^\alpha u\|_{L^\infty} .
\]
有限$p$的范数累计整个区域上的积分，$p=\infty$的范数只看各分量的本质上界。例如$\|\,|x|\,\|_{W^{1,\infty}(-1,1)}=\max(1,1)=1$。由定理~\ref{thm:sob-1d-absolute-continuity}，一维区间上$W^{1,\infty}$的元素恰好是有Lipschitz代表的函数。

\subsection{完备性与Hilbert结构}
\label{sec:sob-completeness}

第~\ref{sec:fs-lp-completeness}节证明了$L^p$完备。Sobolev空间的元素是一组$L^p$分量，各分量之间还被弱导数关系捆在一起；完备性要求极限之后，这种捆绑仍然成立。

\begin{theorem}[Sobolev空间的完备性]
\label{thm:sob-complete}
对任意开集$\Omega\subset\mathbb R^d$、整数$k\ge0$和$1\le p\le\infty$，$W^{k,p}(\Omega)$是Banach空间。
\end{theorem}

\begin{proof}
设$(u_n)$是$W^{k,p}(\Omega)$中的Cauchy列。

第一步，取得分量极限。每个分量的误差不超过全范数，即$\|D^\alpha u_n-D^\alpha u_m\|_{L^p}\le\|u_n-u_m\|_{W^{k,p}}$，所以对每个$|\alpha|\le k$，$(D^\alpha u_n)_n$是$L^p$中的Cauchy列。由$L^p$的完备性，存在$v_\alpha\in L^p(\Omega)$使$D^\alpha u_n\to v_\alpha$；记$v=v_0$为$u_n$自身的极限。此时这些$v_\alpha$只是一组各自取得的$L^p$元素，还不知道它们与$v$有什么关系。

第二步，把弱导数恒等式传到极限。固定测试$\varphi$，对每个$n$有
\[
 \int_\Omega D^\alpha u_n\,\varphi\,\mathrm dx=(-1)^{|\alpha|}\int_\Omega u_n\,D^\alpha\varphi\,\mathrm dx .
\]
设$q$为$p$的共轭指数（$p=1$时$q=\infty$，$p=\infty$时$q=1$）。由Hölder不等式，
\[
 \begin{aligned}
 \Bigl|\int_\Omega(D^\alpha u_n-v_\alpha)\varphi\,\mathrm dx\Bigr|&\le\|D^\alpha u_n-v_\alpha\|_{L^p}\|\varphi\|_{L^q},\\
 \Bigl|\int_\Omega(u_n-v)D^\alpha\varphi\,\mathrm dx\Bigr|&\le\|u_n-v\|_{L^p}\|D^\alpha\varphi\|_{L^q} .
 \end{aligned}
\]
测试$\varphi$光滑且支撑紧，$\varphi$与$D^\alpha\varphi$有界并只在有限测度的集合上非零，所以对每个$1\le q\le\infty$，$\|\varphi\|_{L^q}$和$\|D^\alpha\varphi\|_{L^q}$都有限；$p=1$时用的是它们的上确界，$p=\infty$时用的是它们的积分。两式右端都随$n\to\infty$趋于零，于是两边取极限得
\[
 \int_\Omega v_\alpha\,\varphi\,\mathrm dx=(-1)^{|\alpha|}\int_\Omega v\,D^\alpha\varphi\,\mathrm dx\qquad(\forall\varphi\in C_c^\infty(\Omega)).
\]
这正是说$v_\alpha=D^\alpha v$。分量极限之间的联系由弱导数的定义重新建立，而不是凭记号默认。

第三步，合并分量。所有$D^\alpha v=v_\alpha\in L^p$，所以$v\in W^{k,p}(\Omega)$。有限$p$时，
\[
 \|u_n-v\|_{W^{k,p}}^p=\sum_{|\alpha|\le k}\|D^\alpha u_n-v_\alpha\|_{L^p}^p\longrightarrow0 ,
\]
$p=\infty$时把求和换成有限个量的最大值，结论相同。所以$u_n\to v$在$W^{k,p}$中成立。
\end{proof}

证明中没有用到$\Omega$的边界形状，也没有得到逐点收敛或边界值。这些都不是完备性能给出的结论。

\paragraph{$H^k$的内积}
$p=2$时，Sobolev范数来自内积。记$H^k(\Omega)=W^{k,2}(\Omega)$，定义
\begin{equation}
 \langle u,v\rangle_{H^k}=\sum_{|\alpha|\le k}\int_\Omega D^\alpha u\,\overline{D^\alpha v}\,\mathrm dx ,
 \label{eq:sob-hk-inner-product}
\end{equation}
复值时对第二个变量取共轭，与第~\ref{sec:fs-hilbert}节的约定一致。每一项都是$L^2$内积，因此有限，并对第一个变量线性、共轭对称；零阶项保证$\langle u,u\rangle_{H^k}=0$推出$u=0$。它诱导的范数正是$\|u\|_{W^{k,2}}$。结合定理~\ref{thm:sob-complete}，$H^k(\Omega)$是Hilbert空间。

只把几个$L^2$内积相加，得到的是一个内积空间；称它为Hilbert空间，靠的是上面的完备性证明，即极限中导数关系仍然成立。$H^1$的内积写开来就是$\langle u,v\rangle_{H^1}=\int u\bar v+\int\nabla u\cdot\overline{\nabla v}$，函数项和梯度项同等计入。

\paragraph{Hilbert结构能给出什么}
有了Hilbert结构，第~\ref{sec:fs-projection}节的正交投影和第~\ref{sec:lfo-riesz}节的Riesz表示都可以直接用于$H^k(\Omega)$：例如$H^1$中的闭子空间有正交补，$H^1$上的有界线性泛函由某个$H^1$元素以内积表示。由第~\ref{sec:lfo-weak-subsequence}节，$H^k$中的有界序列有弱收敛子列。

但空间完备、有弱收敛子列，都不自动给出更强的结论。在$(0,1)$上取$u_n(x)=\sin(n\pi x)/(n\pi)$，则$\|u_n'\|_{L^2}=1/\sqrt2$，$\|u_n\|_{L^2}=1/(\sqrt2\,n\pi)$，序列在$H^1$中有界。对任意$v\in H^1(0,1)$，$\langle u_n,v\rangle_{H^1}=\int u_n\bar v+\int\cos(n\pi x)\,\overline{v'(x)}\,\mathrm dx$，第一项被$\|u_n\|_{L^2}\|v\|_{L^2}$控制而趋零，第二项是$\overline{v'}\in L^2$的余弦系数，由Bessel不等式也趋零，所以$u_n\rightharpoonup0$。可是$\|u_n\|_{H^1}\to1/\sqrt2$，序列不在$H^1$范数下收敛。值得注意的是，它在$L^2$中确实收敛到$0$。导数有界的序列能否在较弱的范数下抽出强收敛子列，函数是否有连续代表，边界值是否有意义，这些问题依赖于区域和指数，分别由第~\ref{sec:sob-regularity}节处理。在此之前，我们先建立几乎所有这些论证都要用到的工具：光滑逼近。
